\section*{Цель работы}
Экспериментальное определение параметров двух индуктивно связанных катушек и проверка основных соотношений для цепей с индуктивной связью при различных соединениях катушек.

\section*{Задачи работы}
\begin{enumerate}
    \item Определение собственных индуктивностей катушек ($L_1$, $L_2$), их взаимной индуктивности ($M$) и коэффициента связи ($k$).
    \item Экспериментальное исследование последовательного (согласного и встречного) соединения индуктивно связанных катушек.
    \item Экспериментальное исследование параллельного (согласного и встречного) соединения индуктивно связанных катушек.
    \item Исследование амплитудно-частотной характеристики (АЧХ) трансформатора при работе на активную нагрузку.
\end{enumerate}

\section*{Теоретические сведения}
Индуктивно связанные цепи --- это цепи, в которых изменение тока в одном элементе (первичной обмотке) приводит к появлению ЭДС в другом элементе (вторичной обмотке) за счет явления взаимной индукции. Схема замещения двух индуктивно связанных катушек представлена на \cref{fig:equiv_circuit}.

\begin{figure}[H]
    \centering
    \includegraphics[width=0.7\linewidth]{figs/equiv_circuit.pdf}
    \caption{Схема замещения двух индуктивно связанных катушек}
    \label{fig:equiv_circuit}
\end{figure}

Степень магнитной связи между катушками характеризуется коэффициентом связи $k$:
\begin{equation}
    k = \frac{|M|}{\sqrt{L_1 L_2}} = \frac{|x_M|}{\sqrt{x_1 x_2}}, \quad 0 \le k \le 1
    \label{eq:k_factor}
\end{equation}
где $x_1 = \omega L_1$ и $x_2 = \omega L_2$ --- индуктивные сопротивления катушек, $x_M = \omega M$ --- сопротивление взаимной индуктивности.

В режиме гармонических колебаний система описывается уравнениями:
\begin{equation}
    \begin{cases}
        \dot{U}_1 = (R_1 + j\omega L_1)\dot{I}_1 + j\omega M \dot{I}_2 = (R_1 + jx_1)\dot{I}_1 + jx_M \dot{I}_2 \\
        \dot{U}_2 = j\omega M \dot{I}_1 + (R_2 + j\omega L_2)\dot{I}_2 = jx_M \dot{I}_1 + (R_2 + jx_2)\dot{I}_2
    \end{cases}
    \label{eq:system}
\end{equation}
Знак $M$ (и $x_M$) зависит от способа включения катушек. При согласном включении (токи входят или выходят из однополярных выводов, отмеченных на схеме) $M > 0$. При встречном включении $M < 0$.

Параметры катушек определяют из опытов холостого хода, пренебрегая активными сопротивлениями обмоток ($R_1 \approx 0, R_2 \approx 0$).
\begin{itemize}
    \item \textit{Опыт 1 (вторая катушка разомкнута, $\dot{I}_2=0$):} $x_1 = \frac{U_1}{I_1}$, $|x_M| = \frac{U_2}{I_1}$.
    \item \textit{Опыт 2 (первая катушка разомкнута, $\dot{I}_1=0$):} $x_2 = \frac{U_2}{I_2}$, $|x_M| = \frac{U_1}{I_2}$.
\end{itemize}

При последовательном соединении катушек (\cref{fig:connections}) эквивалентная индуктивность равна:
\begin{equation}
    L_\text{э} = L_1 + L_2 + 2M
    \label{eq:series}
\end{equation}

При параллельном соединении катушек (\cref{fig:connections}) эквивалентная индуктивность равна:
\begin{equation}
    L_\text{э} = \frac{L_1 L_2 - M^2}{L_1 + L_2 - 2M}
    \label{eq:parallel}
\end{equation}
В формулах \cref{eq:series} и \cref{eq:parallel} $M>0$ для согласного включения и $M<0$ для встречного.

\begin{figure}[H]
    \centering
    \includegraphics[width=\linewidth]{figs/connections.pdf}
    \caption{Последовательное (а) и параллельное (б) соединения катушек}
    \label{fig:connections}
\end{figure}

Если ко второй катушке подключить нагрузку $Z_\text{н}$, получается трансформатор (\cref{fig:transformer_circuit}). Его свойства описываются функцией передачи по напряжению $H_U(j\omega) = \dot{U}_2 / \dot{U}_1$. При активной нагрузке $Z_\text{н} = R_\text{н}$ и пренебрежении потерями в катушках модуль функции передачи (АЧХ) имеет вид:
\begin{equation}
    |H_U(j\omega)| = \frac{|M| R_\text{н}}{\sqrt{\omega^2(L_1 L_2 - M^2)^2 + (L_1 R_\text{н})^2}}
    \label{eq:afc}
\end{equation}

\begin{figure}[H]
    \centering
    \includegraphics[width=0.6\linewidth]{figs/transformer_circuit.pdf}
    \caption{Схема трансформатора с нагрузкой}
    \label{fig:transformer_circuit}
\end{figure}

\section{Порядок выполнения работы}
\subsection{Определение параметров катушек}
\begin{enumerate}
    \item Соберите схему, показанную на \cref{fig:exp_setup}. В качестве источника используется генератор сигналов (ГС), измерения проводятся цифровыми вольтметрами и амперметром.
    \begin{figure}[H]
        \centering
        \includegraphics[width=0.7\linewidth]{figs/exp_setup.pdf}
        \caption{Схема для определения параметров катушек}
        \label{fig:exp_setup}
    \end{figure}
    \item Установите на генераторе синусоидальное напряжение с частотой $f = 1~\text{кГц}$.
    \item Подключите выводы А и Б генератора к клеммам $1-1'$ первой катушки, установив выходное напряжение ГС таким образом, чтобы напряжение на катушке было $U_1 \approx 2~\text{В}$. Клеммы $2-2'$ второй катушки оставьте разомкнутыми.
    \item Измерьте напряжение на первой катушке $U_1$, ток в ней $I_1$ и напряжение на второй (разомкнутой) катушке $U_2$. Занесите данные в \cref{tab:params}.
    \item Подключите выводы А и Б генератора к клеммам $2-2'$ второй катушки, установив напряжение $U_2 \approx 2~\text{В}$. Клеммы $1-1'$ первой катушки оставьте разомкнутыми.
    \item Измерьте напряжение на второй катушке $U_2$, ток в ней $I_2$ и напряжение на первой (разомкнутой) катушке $U_1$. Занесите данные в \cref{tab:params}.
\end{enumerate}

\subsection{Исследование последовательного соединения}
\begin{enumerate}
    \item Соберите схему последовательного соединения катушек, как на \cref{fig:connections} (а).
    \item Осуществите \textit{согласное} включение, соединив клемму $1'$ с клеммой $2$. Подайте напряжение от ГС ($f = 1~\text{кГц}$) на клеммы $1$ и $2'$, установив его равным $U = 2~\text{В}$.
    \item Измерьте общее напряжение $U$, общий ток $I$, а также напряжения на каждой катушке $U_1$ и $U_2$. Занесите данные в \cref{tab:series_conn}.
    \item Осуществите \textit{встречное} включение, поменяв местами выводы одной из катушек (например, соединив $1'$ с $2'$).
    \item Повторите измерения из п. 3 и занесите данные в \cref{tab:series_conn}.
\end{enumerate}

\subsection{Исследование параллельного соединения}
\begin{enumerate}
    \item Соберите схему параллельного соединения катушек, как на \cref{fig:connections} (б).
    \item Осуществите \textit{согласное} включение, соединив клеммы $1$ и $2$, а также $1'$ и $2'$. Подайте напряжение от ГС ($f = 1~\text{кГц}$) на соединенные клеммы, установив его равным $U = 1~\text{В}$.
    \item Измерьте общее напряжение $U$ и общий ток в цепи $I$. Занесите данные в \cref{tab:parallel_conn}.
    \item Осуществите \textit{встречное} включение (например, соединив $1$ с $2'$, а $1'$ с $2$).
    \item Повторите измерения из п. 3 и занесите данные в \cref{tab:parallel_conn}.
\end{enumerate}

\subsection{Исследование АЧХ трансформатора}
\begin{enumerate}
    \item Соберите схему трансформатора (\cref{fig:transformer_circuit}). Первичную обмотку ($1-1'$) подключите к ГС.
    \item В качестве нагрузки $R_\text{н}$ ко вторичной обмотке ($2-2'$) подключите резистор сопротивлением $R_{\text{н}1} = 100~\text{Ом}$.
    \item Установите на выходе ГС напряжение $U_1 = 1~\text{В}$.
    \item Изменяя частоту генератора в диапазоне от $100~\text{Гц}$ до $10~\text{кГц}$, измеряйте напряжение на вторичной обмотке $U_2$. В процессе измерений поддерживайте напряжение на первичной обмотке постоянным ($U_1 = 1~\text{В}$). Результаты занесите в \cref{tab:afc}.
    \item Замените нагрузочный резистор на $R_{\text{н}2} = 1000~\text{Ом}$ и повторите измерения из п. 4. Результаты занесите в \cref{tab:afc}.
\end{enumerate}

\section{Обработка результатов}
\begin{enumerate}
    \item По данным из \cref{tab:params} вычислите индуктивные сопротивления $x_1$, $x_2$, два значения $|x_M|$, индуктивности $L_1$, $L_2$, два значения $|M|$. Найдите среднее значение $|M|_\text{ср}$.
    \item Используя среднее значение $|M|_\text{ср}$, рассчитайте коэффициент связи $k$ по \cref{eq:k_factor}.
    \item По данным из \cref{tab:series_conn} рассчитайте экспериментальные значения эквивалентной индуктивности $L_\text{э} = U / (\omega I)$ для согласного и встречного включений. Сравните их с теоретическими значениями, рассчитанными по \cref{eq:series}.
    \item По данным из \cref{tab:parallel_conn} рассчитайте экспериментальные значения эквивалентной индуктивности $L_\text{э} = U / (\omega I)$ для согласного и встречного включений. Сравните их с теоретическими значениями, рассчитанными по \cref{eq:parallel}.
    \item По данным из \cref{tab:afc} для каждой частоты и каждой нагрузки рассчитайте экспериментальное значение модуля коэффициента передачи $|H_U(j\omega)|_\text{эксп} = U_2 / U_1$.
    \item Для тех же частот и нагрузок рассчитайте теоретические значения $|H_U(j\omega)|_\text{теор}$ по \cref{eq:afc}, используя ранее найденные параметры $L_1, L_2, M$.
    \item Постройте на двух разных графиках (для $R_{\text{н}1}$ и $R_{\text{н}2}$) зависимости $|H_U(j\omega)|_\text{эксп}$ и $|H_U(j\omega)|_\text{теор}$ от частоты $f$.
    \item Сформулируйте выводы по работе и ответьте на контрольные вопросы из методического пособия.
\end{enumerate}

% Соответствие изображений из методического пособия и label в документе:
% Рис. 9.1 -> \label{fig:equiv_circuit}
% Рис. 9.2 -> \label{fig:connections}
% Рис. 9.3 -> \label{fig:transformer_circuit}
% Рис. 9.4 -> \label{fig:exp_setup}